\chapter{Projects and scientific workflows}\label{ch:workflows}
This chapter describes additions in the current source build. Rebuild the application and install the matching Python source package before using them; previously downloaded release binaries do not acquire new features automatically.

\section{Save a project and recover a session}
Use the File menu's project commands to save all open tabs in one \code{.afproject} file. A project stores atomic structures, grain colours and region IDs, camera framing, basic display switches, electronic primary and reference fields, the current mesh and result table, slice settings, and electronic operation history. Opening a project creates tabs without replacing unrelated open work. Projects store the numerical fields themselves, so the original CHGCAR or Cube file is not required to restore the saved calculation.

Use the electronic workspace's Open/Save workspace buttons to exchange a single analysis independently. A project is not a copy of every application preference: per-element material settings, every builder's transient inputs, and running calculations are not serialized. Save important builder parameters alongside the structure. A calculation must finish or be cancelled before a project can be saved. File writes replace the destination only after successful serialization; a failed write preserves the previous project.

Enable autosave in the File menu to write a recovery project approximately once per minute while calculations are idle. The recovery command loads the latest completed snapshot. On Windows the recovery file is below \code{LOCALAPPDATA/AtomForge}; on Linux it is below \code{XDG_STATE_HOME/AtomForge}, falling back to \code{HOME/.local/state/AtomForge}. Autosave is a recovery aid, not a versioned backup system. Work after the latest successful snapshot may need to be repeated.

\section{Electronic history, cancellation and export}
Electronic Undo and Redo restore the previous numerical workspace and views. Interactive history keeps at most four snapshots and clears the undo stack for a state larger than 32 MiB; save a project before expensive operations on large grids. The history records operation selections, numeric parameters and input filenames. It does not independently establish the provenance or physical validity of an external file.

Cancel calculation requests cooperative cancellation. The current stage's progress is shown where the algorithm can measure it. The completed result is discarded if cancellation was requested. Loading and long grid loops check cancellation, although external blocking I/O and uninstrumented work may delay the response.

The Appearance quality selector offers four resolutions. It changes the displayed volume texture, ray-sampling budget and 2D sampling grid; the underlying numerical field and numerical exports retain their original resolution. Use lower quality for interaction and higher quality for a final view. Rendering quality cannot recover detail absent from the input file.

The export selector includes a 3D PNG and a 2D SVG. The PNG is rendered at 1600 by 1200 pixels using the current camera, palette and opacity. It exports the data viewport rather than the whole application dialog. The SVG preserves the current section, pan, zoom, rotation and contour display. For a rotatable downstream model export OBJ or PLY instead. Save a project to retain all loaded and derived fields together.

\section{Structural analysis and automation}
Common-neighbour and angular-distribution analysis are available from Analysis. Interface Builder and Stacking Faults are available from Build. Their C++ algorithms are also exposed through native CLI and Python workflows.
\begin{lstlisting}[language=Python]
import atomforge as af
host = af.load("POSCAR")
rows = af.cna(host, cutoff=3.1)
rdf = af.rdf(host, cutoff=6, bins=120)
adf = af.adf(host, cutoff=3.1, centre="Cu")
sro = af.sro(host, shells=3, tolerance=0.1)
voids = af.interstitial_sites(host, resolution=14, clearance=0.6)
cut = af.sculpt(host, [(1,0,0,-5,5), (0,1,0,-5,5),
                      (0,0,1,-5,5)])
\end{lstlisting}
Cutoff lengths and sculpt bounds are in Angstrom. Manual sculpt bounds are relative to the supercell centroid; periodic bounds instead specify an integer start plane and a number of structural periods. CNA is cutoff-sensitive: include the complete intended neighbour shell and use a sufficiently large periodic cell. FCC/HCP and BCC require different neighbour-shell conventions. ADF counts unordered neighbour pairs; smoothing affects the displayed values, not raw counts. These geometrical analyses do not replace energetic validation.

The native \code{--analyze} modes are cna, rdf, adf, sro, interstitial and sculpt. All require \code{--input} and \code{--output}; analyses write CSV, while sculpt writes a structure. Unsupported options are rejected. Native \code{--convert} provides structure conversion and extended XYZ export. The generated API appendix lists Python parameters and all ten builder modes, including interface and stacking-fault sequence export.

\section{Recipes and batch electronic processing}
\begin{lstlisting}[language=Python]
from atomforge.electronic import load_volume, run_pipeline, batch_process
density = load_volume("CHGCAR").fields[0]
steps = [
    {"operation":"scale", "parameters":{"factor":2}, "name":"twice"},
    {"operation":"subtract", "reference":"input", "name":"difference"},
    {"operation":"integrate", "name":"electrons"},
]
results = run_pipeline(density, steps)
report = batch_process(["density1.cube", "density2.cube"],
                       steps, "batch-results")
\end{lstlisting}
Every step names an allowlisted operation, optional parameters and an optional unique result name. A reference can identify an earlier grid or an external file. The original input remains available as \code{input}; intermediate results can be inspected before saving. Geometry and units must match for binary operations; resampling is explicit. Batch outputs use input stems, reject name collisions and existing destinations, and report each input's success or failure. Inputs named CHGCAR in different folders need distinct output names or separate batch destinations.

The Python electronic command accepts an operation, a reference, JSON parameters, a JSON recipe or batch inputs. Use its \code{--help} for current argument syntax. Recipe execution does not evaluate Python expressions. Save the recipe, source calculation settings and input hashes separately when preparing an independently reproducible study.

\section{Trajectories and animation}
Open Analysis > Trajectories and transport > Trajectory playback, choose an XYZ or extended XYZ trajectory, and use Play/Pause, First frame, the frame slider and playback speed. Frames replace the active structure; preserve the original in another tab or project first. Main-window orbit, zoom and pan remain available. The desktop reader requires species followed by Cartesian coordinates; extended XYZ properties must begin with \code{species:S:1:pos:R:3}. It reads lattice vectors from the comment and rejects incomplete frames.

Python provides broader ASE trajectory formats, including XDATCAR and LAMMPS dumps:
\begin{lstlisting}[language=Python]
from atomforge.science import read_trajectory, write_trajectory, export_animation
frames = read_trajectory("XDATCAR", format="vasp-xdatcar")
write_trajectory(frames, "trajectory.extxyz")
export_animation(frames, "trajectory.gif", fps=12, radii=0.5)
\end{lstlisting}
GIF export uses fixed projection bounds across frames and explicitly chosen radii. Conversion preserves species, positions and a full cell, but the current Structure model does not retain per-axis periodicity, velocities or arbitrary per-atom trajectory properties. Keep the original trajectory for dynamics analysis.

\section{Install scientific integrations}
The following workflows live in \code{atomforge.science}; they are Python integrations rather than additional native desktop dialogs. Install the source package with its \code{science} extra to obtain ASE, pymatgen, Matplotlib and Gemmi. The native electronic library is additionally needed for stockholder integration and reciprocal-surface triangulation. External DFT, Bader and Chargemol executables and licensed calculation inputs are not bundled.

Gaussian input and Quantum ESPRESSO input conversion use ASE. A PWI export is a structural template with element-named UPF placeholders; supply real pseudopotentials and calculation parameters before running it. MOL, MOL2 and SDF conversion uses the native executable. CIFs carrying nontrivial space-group information require the crystallography extra (Gemmi) and expand symmetry-related sites. Validate disorder/partial-occupancy models separately because the Structure model represents explicit atoms.

\section{Bands, DOS, projected data and orbitals}
\begin{lstlisting}[language=Python]
from atomforge.science import read_bands, plot_bands, read_dos, plot_dos
bands = read_bands("EIGENVAL")
plot_bands(bands, "bands.svg", fermi_eV=5.2)
dos = read_dos("DOSCAR")
plot_dos(dos, "dos.png")
\end{lstlisting}
Band arrays are indexed by k point then band, separated by spin channel. Without an explicit reciprocal-path distance array, the plot labels its horizontal axis as k-point index. DOSCAR parsing checks finite data, complete blocks and consistent energy grids. Projected DOS column meanings depend on the producer's spin and LORBIT settings; pass \code{projected_labels} when these are known rather than inferring orbital labels from an ambiguous column count.

\code{read_projections("PROCAR")} returns pymatgen's projection object. Select and sum the intended ion/orbital channels to construct nonnegative weights matching the band arrays, then pass these to \code{plot_bands(..., projections=weights)} for fat bands. Spin-orbit projections require the conventions of the supplied PROCAR. Compare the number and ordering of k points and bands before combining files.

\code{orbital_density(wavecar,poscar,output,kpoint=0,band=0)} reconstructs a selected orbital and writes PARCHG for the desktop electronic viewer. Indices are zero-based; spin, spinor and FFT scale are explicit options. This is a plane-wave reconstruction through pymatgen and excludes PAW augmentation; it is not an all-electron density or an exact total-energy calculation.

\section{Reciprocal-space surfaces}
\begin{lstlisting}[language=Python]
from atomforge.science import read_bandgrid, plot_surface
bands = read_bandgrid("bands.bxsf", energy_unit="eV", reciprocal_unit="1/A")
label = next(iter(bands.bands))
plot_surface(bands, label, "fermi-surface.png")
bands.surface(label).save("fermi-surface.obj")
\end{lstlisting}
BXSF BANDGRID data are read in z-fastest order and preserve endpoint samples. The default isosurface level is the file's Fermi energy; out-of-range levels are rejected. Supply the producing program's actual units. The resulting surface spans the reciprocal parallelepiped and is not folded into the first Brillouin zone. Keep it separate from real-space charge grids. Raw BXSF input is intentionally not passed to the desktop real-space density loader.

\section{Charge partitioning}
\code{hirshfeld(density, proatoms)} applies stockholder weights to an electron-density grid using caller-supplied neutral-proatom grids. All geometries and density units must agree. Proatoms must cover every nonzero density sample. The electron populations sum to the integrated density; optional reference electron counts produce net atomic charges. Use compatible valence or all-electron conventions throughout. AtomForge does not invent atomic reference densities from Gaussian functions.

\code{bader(path, reference=..., potcar=..., executable=...)} runs the actual Henkelman solver through pymatgen. A summed AECCAR0+AECCAR2 reference can define basins while CHGCAR supplies the integrated valence charge. Returned raw populations count electrons; when POTCAR is provided, net charge is reference valence count minus population. Inspect vacuum charge and total electron count as convergence checks. Cube input requires the explicit cube format option.

\code{ddec(directory, atomic_densities=...)} uses Chargemol DDEC6/CM5. It requires the actual solver, atomic-density database and required VASP files. Set \code{run=False} to parse existing results. Errors are reported directly; neither Bader nor DDEC falls back to a Voronoi partition. Backend availability and numerical convergence are separate from successful import of the Python wrapper.

\section{Diffraction, relaxation, dynamics and phonons}
\code{powder_diffraction(structure)} uses pymatgen's X-ray or neutron calculator with explicit wavelength and angle limits. The result contains peak angles, intensities and reflection information; \code{plot_diffraction} produces a broadened pattern. \code{single_crystal} uses explicitly supplied per-site, per-reflection complex scattering factors. Include any intended occupancy and displacement attenuation in those effective factors. Radiation, isotope and form-factor assumptions affect intensities and should be recorded.

\begin{lstlisting}[language=Python]
from ase.calculators.emt import EMT
from atomforge.science import relax, molecular_dynamics, phonons, mode_frames
result = relax(host, EMT(), fmax=0.02, steps=200)
print(result.converged, result.energy)
motion = molecular_dynamics(host, EMT(), steps=100,
                            timestep_fs=1, temperature_K=300, seed=7)
modes = phonons(host, EMT(), [[0,0,0]], directory="fresh-phonon-cache",
                supercell=(2,2,2), displacement=0.01)
animation = mode_frames(host, modes["modes"][0,0])
\end{lstlisting}
EMT is an example potential with limited element support, not a general DFT model. A caller-configured ASE calculator can instead connect an external DFT or interatomic-potential engine. Its executable, pseudopotentials, settings and convergence remain the caller's responsibility. BFGS returns a convergence flag; reaching the step limit does not establish equilibrium. Cell relaxation requires a periodic structure and stress-capable calculator.

Dynamics uses NVE velocity-Verlet. Temperature sets the initial velocity distribution, not a thermostat. Check total-energy drift and timestep convergence. Phonons use finite displacements in a fresh cache directory; q points are fractional reciprocal coordinates, energies are in eV, and negative values denote imaginary frequencies. Check supercell and displacement convergence. \code{mode_frames} animates a supplied Gamma-point mode with an explicit maximum displacement amplitude. A geometric builder or a rendered vibration does not by itself establish dynamical stability.

\section{Twenty condensed-matter tools: desktop, Python and command line}
The Analysis menu groups tools by task: \textbf{Structure and defects}, \textbf{Trajectories and transport}, \textbf{Electronic structure}, \textbf{Reciprocal space}, \textbf{Thermal and mechanical}, and \textbf{Simulation}. Existing tools are alongside related analyses: CNA, RDF, ADF and SRO appear with structural analyses; trajectory playback appears with transport analyses; electronic post-processing, Wannier and LOBSTER appear with electronic analyses. The twenty tools below share numerical implementations with \code{atomforge.science}. The desktop launches a separate Python process; select its interpreter on the \textbf{Environment} tab. Install the matching source package with its \code{science} extra into that environment. Python 3.12 is tested. NumPy, SciPy, ASE and SeeK-path provide numerical and simulation support; a Python interpreter without these dependencies cannot run the tools.

Choose a tool from its task section. Its named dialog has \textbf{Inputs}, \textbf{Results}, \textbf{Method} and \textbf{Environment} tabs. Read the Method tab for units and assumptions, then fill the grouped Inputs form. Optional inputs expand when their checkbox is enabled; Boolean settings use one checkbox. Periodic boundaries have separate a, b and c checkboxes. Fit and vacuum intervals use Start/End fields; the spectrum window uses a choice list. Select \textbf{Load from file} and \textbf{Browse} for numerical CSV, NPY, JSON or structure/trajectory files, or select \textbf{Enter values} for small JSON arrays. Expand \textbf{File options} to select a zero-based column or a dotted JSON field such as \code{result.msd_A2}. Numeric CSV files have no header; matrix inputs use the whole table. Only data are read; array strings are not evaluated as Python code. The scrollable form keeps Run calculation accessible on smaller windows, and completed calculations select the Results tab automatically.

Structure inputs in simulation tools offer \textbf{Use active structure}. The application copies the current structure into the calculation request without modifying the active tab. For array-based local analyses, load a structure file for its Cartesian positions and, if periodic, supply the cell and periodic-axis flags separately. Structure files with multiple frames supply trajectory positions. Atom ordering must remain consistent. Velocity analysis requires explicit velocity data in Angstrom/fs or an ASE-readable trajectory containing momenta; positions alone do not supply velocities.

Compressed EIGENVAL files (gzip, bzip2 and xz) are supported. XDATCAR import accepts repeated cell headers. For trajectory files, velocity, mass and cell inputs select the corresponding property automatically. Numeric column selection works for CSV, NPY and JSON tables; invalid columns are rejected. When a trajectory contains \code{time_fs} metadata, its intervals must be uniform and match the requested sampling interval. A simulation's final partial interval must be trimmed or resampled before time-correlation analysis. Simulation structure conversion rejects mixed periodic/nonperiodic axes instead of silently imposing three-dimensional periodicity; nonperiodic bounding boxes remain nonperiodic.

Press \textbf{Run calculation}. Cancel terminates the calculation process and its children. The result panel displays scalar results, small tensors and summary ranges for large arrays; \textbf{Save results} uses the shared file browser to export full JSON arrays, validity masks, structural frames and the original request parameters. For NEB, NVT, NPT and reciprocal paths, \textbf{Save structures} exports extended XYZ and \textbf{Open final structure in a new tab} brings the last frame or primitive structure into the viewport. Saved trajectories can be opened in Trajectory playback. Invalid local environments are represented by null in JSON, with a separate validity mask. Large arrays are not fully printed in the summary. These forms do not currently colour the main viewport by derived per-atom properties or automatically plot all returned curves. Results can be plotted through Python or imported into a plotting package. Temporary calculation files are retained under the system temporary directory's \code{AtomForge-science} folder; save results needed for a permanent record.

\subsection{Trajectory and local-structure analyses (priorities 1--9)}
\begin{enumerate}
\item \textbf{Mean-square displacement (MSD).} Load positions shaped as frames by atoms by three coordinates, and set the sampling interval in fs. Prefer unwrapped coordinates. If wrapped is enabled, supply a fixed cell; the inter-frame motion must be small enough to resolve the shortest periodic displacement. Run and inspect the lag-time MSD and Cartesian displacement tensor. Optional geometric-centroid drift removal is not mass-weighted centre-of-mass removal. A changing cell requires a physically chosen reference-frame transformation before this analysis.
\item \textbf{Diffusion coefficient.} Supply the MSD lag times and values, then explicitly select a diffusive fit interval. Set dimensions to the number of directions summed in that MSD. The result uses $D=\mathrm{slope}/(2d)$ and reports both Angstrom$^2$/fs and m$^2$/s, intercept and $R^2$. Check several fit intervals and independent trajectories. The ordinary least-squares standard error is a fit diagnostic, not a confidence interval accounting for correlated time origins.
\item \textbf{Velocity autocorrelation.} Supply velocities in Angstrom/fs and the frame interval. The result averages $\mathbf v(t)\cdot\mathbf v(t+\tau)$ over atoms and time origins. Optional normalization sets the zero-lag value to one; zero-velocity data cannot be normalized. Truncate the lag range before poorly sampled tails dominate interpretation.
\item \textbf{Trajectory vibrational spectrum.} Load velocities, choose a Hann or rectangular window, and optionally supply atomic masses. The one-sided velocity periodogram is normalized to unit integrated spectral weight per THz. Inspect its reported resolution and Nyquist frequency. This classical velocity spectrum is not an IR, Raman or quantum neutron-scattering intensity.
\item \textbf{Local strain and $D^2_{\min}$.} Load corresponding reference and current coordinates, choose a reference neighbor cutoff, and supply cells and periodic axes. The tool fits the local deformation gradient and reports Green--Lagrange strain and the sum of squared non-affine residuals. Known box deformation is retained. Three-dimensional rank-deficient neighborhoods are marked invalid; a planar neighborhood does not define a full 3D strain tensor.
\item \textbf{Centrosymmetry parameter.} Choose an even neighbor count, commonly twelve for FCC or eight for BCC, and a cutoff containing those neighbors. The conventional Kelchner construction sums the smallest $N/2$ squared pair-vector sums among the nearest $N$ neighbors. Perfect centrosymmetric environments approach zero. Missing neighbors are invalid, not zero; a nonzero value alone does not identify a unique defect type.
\item \textbf{Steinhardt bond order.} Select angular degrees, typically four and six, a cutoff, cell and periodic axes. The tool returns rotationally invariant local $q_l$ from equal-weight bond spherical harmonics. For an ideal FCC first-neighbor shell, $q_4\simeq0.190941$ and $q_6\simeq0.574524$. These are local $q_l$, not neighbor-averaged $\bar q_l$ or third-order $w_l$ invariants.
\item \textbf{Wigner--Seitz defects.} Load reference sites and current atom positions. Supply the reference cell and, if needed, a deformed current cell to remove homogeneous box strain. Inspect site occupancies, vacant sites and sites with more than one occupant. An occupancy $n>1$ contributes $n-1$ excess interstitial atoms. Equidistant assignments are flagged; the tool does not uniquely label which atom in a multiply occupied site is the interstitial.
\item \textbf{Static structure factor.} Supply one structure or a trajectory and Cartesian wave vectors in radians/Angstrom, including the reciprocal $2\pi$ factor. Optional scattering amplitudes weight atoms. The normalization is $S(\mathbf q)=\langle|\sum_jb_j e^{i\mathbf q\cdot\mathbf r_j}|^2\rangle/\sum_j|b_j|^2$; equal weights give $S(0)=N$. Use commensurate reciprocal points for wrapped periodic coordinates. No powder averaging or forward-scattering subtraction is implied.
\end{enumerate}

\subsection{Electronic and thermomechanical analyses (priorities 10--16)}
\begin{enumerate}\setcounter{enumi}{9}
\item \textbf{Band gap and band edges.} Load an EIGENVAL file or an energy array shaped as k points by bands, optionally preceded by a spin dimension. Supply the Fermi energy separately in the same energy reference. Inspect the fundamental sampled gap, smallest same-spin direct gap and extremum indices. Band crossings and states at the Fermi tolerance are reported as metallic. A sparse k mesh can miss pockets or extrema, and the direct gap is not an optical excitation energy.
\item \textbf{Effective-mass tensor.} Supply at least ten independent points in a small three-dimensional Cartesian k neighborhood, one consistently tracked band's energies, and the expansion centre. Inspect signed principal masses in electron-mass units, principal axes, the residual band gradient and fit error. A one-dimensional plotted band path is insufficient for the full tensor. Negative electron curvature corresponds to positive hole mass after sign reversal.
\item \textbf{Work function.} First obtain a planar electron potential-energy profile, for example using the electronic workspace's planar averaging. Supply distances, potential values, Fermi energy and a vacuum interval. The result is the mean vacuum level minus the Fermi energy. Inspect the slope/flatness flag and vacuum spread. Evaluate opposite slab faces separately; dipole corrections are not introduced by this analysis.
\item \textbf{Equation of state.} Supply at least five distinct positive volumes and consistent total energies bracketing a minimum. A third-order Birch--Murnaghan fit reports equilibrium volume, energy, bulk modulus in GPa, pressure derivative and residuals. Compare residuals and vary the fitted volume range. Energies and volumes must refer to the same cell or atom-count normalization.
\item \textbf{Elastic tensor and moduli.} Supply independently varied small strains and the corresponding stresses in GPa. Engineering-Voigt order is xx, yy, zz, yz, xz, xy; shear strains are twice tensor shear, while shear stresses are not doubled. Stress is positive in tension. The fit includes prestress, symmetrizes stiffness, and reports mismatch, residuals, zero-prestress stability and Voigt--Reuss--Hill moduli where stable. Check strain-amplitude convergence; the stability test does not substitute for finite-pressure stability criteria.
\item \textbf{Phonon density of states.} Supply mode energies in eV shaped as q points by branches, a plotting grid, Gaussian width and optional q-point weights. The analytic total weight is the number of branches per cell. Inspect the enclosed plotted weight to detect truncated tails. Negative imaginary-mode energies remain on the negative axis; broadening does not make them stable.
\item \textbf{Harmonic thermodynamics.} Supply a converged mode mesh, normalized or unnormalized nonnegative q weights and a temperature list. Results include zero-point energy, Helmholtz free energy, internal energy, entropy and constant-volume heat capacity, per cell. At zero temperature, free/internal energies equal the zero-point energy. Imaginary modes beyond tolerance are rejected; omitted near-zero mode weight is reported. Electronic, configurational, thermal-expansion and anharmonic contributions are absent.
\end{enumerate}

\subsection{Simulation and reciprocal-space workflows (priorities 17--20)}
\begin{enumerate}\setcounter{enumi}{16}
\item \textbf{NEB migration path.} Select relaxed initial and final structures with the same ordered species and fixed cell. Choose a calculator factory, image count, spring constant, force tolerance and iteration limit. The tool performs improved-tangent NEB with optional climbing image and shortest-image interpolation. Inspect convergence, endpoint forces and both sampled barriers. Each image gets a separate calculator; configured calculator directories are separated into image subdirectories by the workflow adapter. External files, pseudopotentials and endpoint relaxation remain the user's responsibility.
\item \textbf{NVT dynamics.} Select a structure and explicit ASE calculator, temperature, timestep and thermostat relaxation time. Langevin dynamics returns frames, physical energies, temperatures and velocities. The seed controls reproducibility of initialization and thermostat noise. Equilibrate before collecting statistics; thermostatting does not establish equilibration or conserve the physical energy.
\item \textbf{NPT dynamics.} Select a fully periodic structure and a stress-capable calculator, temperature, pressure in GPa and thermostat/barostat times. The isotropic Martyna--Tobias--Klein integrator changes cell volume while keeping shape fixed. Positive target pressure is compression. Inspect temperature, pressure and volume histories and convergence with timestep; instantaneous fluctuations should not be mistaken for failure to match the target.
\item \textbf{Symmetry reciprocal-space path.} Select a periodic structure, a reciprocal sampling spacing and symmetry tolerance. SeeK-path generates an HPKOT path and returns the standardized primitive structure, reciprocal basis, labels and explicit disconnected segments. Use its returned primitive cell when interpreting fractional k points; input supercell fractions are generally different. The time-reversal switch controls path completion, not magnetic-space-group analysis.
\end{enumerate}

\subsection{Repeatable commands and Python examples}
All twenty functions are exported from \code{atomforge.science}. The catalog command prints their parameter names, defaults and physical assumptions. A request is a JSON object with those parameter names; data references are relative to the request file. Example \code{gap-request.json}:
\begin{lstlisting}
{"energies_eV": {"file": "EIGENVAL"}, "fermi_eV": 5.4}
\end{lstlisting}
\begin{lstlisting}[language=bash]
python -m atomforge.science --catalog
python -m atomforge.science band-gap --input gap-request.json --output gap.json
\end{lstlisting}
Existing outputs require an explicit \code{--overwrite}. A failed calculation preserves the previous result. Result JSON stores the tool identifier, parameters and full result. A data reference can select a column or a nested result field:
\begin{lstlisting}
{"file":"profile.csv", "column":1}
{"file":"msd.json", "field":"result.msd_A2"}
\end{lstlisting}
Numerical CSV has no header; NPY loading forbids pickled Python objects.
\begin{lstlisting}[language=Python]
from atomforge.science import mean_square_displacement, diffusion_coefficient
from atomforge.science.workflows import run_tool
msd = mean_square_displacement(unwrapped_positions, timestep_fs=2,
                               max_lag=200)
fit = diffusion_coefficient(msd["lag_fs"], msd["msd_A2"],
                            fit_range_fs=(100, 300), dimensions=3)
print(fit["D_m2_per_s"], fit["r_squared"])
gap = run_tool("band-gap", {"energies_eV": {"file": "EIGENVAL"},
                            "fermi_eV": 5.4})
\end{lstlisting}
The direct Python simulation functions accept configured ASE calculators. Desktop/CLI requests specify their factory, for example:
\begin{lstlisting}
{"module":"ase.calculators.emt", "attribute":"EMT"}
\end{lstlisting}
Optional keyword arguments go under \code{kwargs}. This is explicit calculator selection, not a universal potential. EMT has limited element support. For NEB the direct Python interface takes a callable that returns a new calculator for each image.
